\documentclass[../../../include/open-logic-section]{subfiles}

\begin{document}

\olfileid{sth}{replacement}{strength}
\olsection{પ્રતિસ્થાપનની શક્તિ}

પ્રતિસ્થાપનની શક્તિ વિશેની એક સરળ નોંધથી શરૂઆત કરીએ:
$\Z$થી આગળ ન જઈએ, તો $\omega + \omega$ કરતાં મોટી કે
તેના જેટલી કોઈ પણ વૉન નોયમન ક્રમવાચક સંખ્યાનું અસ્તિત્વ
સાબિત કરી શકતા નથી.

તેનું કારણ રૂપરેખામાં જોઈએ. ~$\ZF$માં કામ કરીને
$V_{\omega+\omega}$ ગણ લો. આ ગણ~$\Z$ના એક \emph{નિદર્શ}
માટે પ્રદેશનું કામ કરે છે. આ સમજવા સૂત્રના
\emph{સાપેક્ષીકરણ}નો સંકેત રજૂ કરીએ:
\begin{defn}\ollabel{formularelativization}
	કોઈ પણ ગણ $M$ અને સૂત્ર $\phi$ માટે, $\phi$ના બધા
	પરિમાણકોને $M$ સુધી મર્યાદિત કરતાં મળતા સૂત્રને $\phi^M$
	કહીએ. એટલે ``$\lexists[x]$''ને ``$(\lexists[x \in M])$''થી
	અને ``$\lforall[x]$''ને ``$(\lforall[x \in M])$''થી બદલો.
\end{defn}\noindent
બતાવી શકાય છે કે~$\Z$ના દરેક સ્વયંસિદ્ધિ $\phi$ માટે
$\ZF \vdash
\phi^{V_{\omega+\omega}}$. પરંતુ
\olref[spine][rank]{ordsetrankalpha} મુજબ $\omega+\omega$
એ $V_{\omega+\omega}$માં \emph{સભ્ય નથી}. તેથી
$\omega+\omega$નું અસ્તિત્વ ન હોવું $\Z$ સાથે સુસંગત છે.

આથી જ \olref[ordinals][replacement]{sec}માં કહ્યું હતું કે
\olref[ordinals][ordtype]{thmOrdinalRepresentation}
પ્રતિસ્થાપન વિના સાબિત કરી શકાતું નથી. કારણ કે~$\Z$માં
સહજ રીતે ક્રમપ્રકાર $\omega+\omega$ \emph{હોવો જોઈએ} એવો
સ્પષ્ટ સુક્રમ વ્યાખ્યાયિત કરવો સરળ છે. ખરેખર,
\olref[ordinals][idea]{sec}માં પ્રાકૃતિક સંખ્યાઓ પરનો આ
ક્રમ આપીને તેનું અનૌપચારિક ઉદાહરણ આપ્યું હતું:
\begin{align*}
	n \lessdot m \text{ જો અને માત્ર જો }&\text{કાં તો }n < m\text{ અને }m-n\text{ યુગ્મ હોય,}\\
	& \text{અથવા $n$ યુગ્મ હોય અને $m$ અયુગ્મ હોય.}
\end{align*}
પરંતુ $\omega+\omega$ અસ્તિત્વમાં ન હોય, તો આ સુક્રમ કોઈ
પણ ક્રમવાચક સંખ્યા સાથે ક્રમ-એકરૂપ નથી. તેથી $\Z$માં
\olref[ordinals][ordtype]{thmOrdinalRepresentation}
\emph{સાબિત થતું નથી}.

ઊલટેથી જોઈએ: પ્રતિસ્થાપનથી $\omega+\omega$નું અસ્તિત્વ
સાબિત થાય છે, અને તેથી $V_{\omega+\omega}$નું અસ્તિત્વ પણ
સાબિત થવું જોઈએ. એટલું જ નહીં. આપણે વ્યાખ્યાયિત કરી શકીએ
તેવા \emph{કોઈ પણ} સુક્રમ માટે
\olref[ordinals][ordtype]{thmOrdinalRepresentation} કહે છે કે
તેની સાથે ક્રમ-એકરૂપ કોઈ $\alpha$ છે, અને તેથી $V_\alpha$
અસ્તિત્વમાં છે. આમ સીધી રીતે, પ્રતિસ્થાપન ખાતરી આપે છે કે
ગણોનો પદાનુક્રમ \emph{ઘણો ઊંચો} હોવો જોઈએ.

આગળના થોડા વિભાગોમાં અને પછી ફરી
\olref[card-arithmetic][fix]{sec}માં પ્રતિસ્થાપન પદાનુક્રમને
\emph{કેટલો} ઊંચો બનાવે છે તે વધુ સમજાશે. હાલનો સરળ મુદ્દો
એ છે કે પ્રતિસ્થાપન માટે ઔચિત્ય \emph{ખરેખર} જોઈએ છે!

\end{document}
